\section*{Chapter \ref{ch:b1:diffeq} --- Linear Differential Equations}

\begin{solution}{exo:b1:diffeq:1}
Homogeneous: $y_h = \lambda\,\eu^{-2x}$. Particular: try $y_p =
c\,\eu^{3x}$ ($3$ is not a root of $r + 2$): $3c + 2c = 1$, $c =
\frac15$. General solution: $y = \frac{\eu^{3x}}{5} +
\lambda\,\eu^{-2x}$. With $y(0) = 1$: $\frac15 + \lambda = 1$,
$\lambda = \frac45$, so $y = \frac{\eu^{3x} + 4\,\eu^{-2x}}{5}$.
\end{solution}

\begin{solution}{exo:b1:diffeq:2}
On $\intoo{0}{+\infty}$, divide by $x$: $y' - \frac{2}{x}\,y = x^2$.
Here $A(x) = -2\ln x$, $\eu^{-A(x)} = x^2$: homogeneous solutions
$\lambda x^2$. Variation of constants: $\mu'(x) = x^2 \cdot x^{-2} =
1$, so $\mu = x + \lambda$ and
\[
y(x) = x^3 + \lambda x^2, \qquad \lambda \in \R .
\]
\emph{Check:} $x(3x^2 + 2\lambda x) - 2(x^3 + \lambda x^2) = x^3$.
\end{solution}

\begin{solution}{exo:b1:diffeq:3}
$y'' - 3y' + 2y = 0$: roots $1$ and $2$; $\;y = \lambda\,\eu^{x} +
\mu\,\eu^{2x}$.

$y'' + 4y' + 4y = 0$: double root $-2$; $\;y = (\lambda + \mu
x)\,\eu^{-2x}$.

$y'' - 2y' + 5y = 0$: roots $1 \pm 2\iu$; $\;y = \eu^{x}(\lambda\cos
2x + \mu\sin 2x)$.
\end{solution}

\begin{solution}{exo:b1:diffeq:4}
Homogeneous: roots $\pm 1$, $y_h = \lambda\,\eu^x + \mu\,\eu^{-x}$.
Particular with polynomial right-hand side ($\gamma = 0$ not a root):
$y_p = ax^2 + bx + c$; substituting, $2a - (ax^2 + bx + c) = x^2$
gives $a = -1$, $b = 0$, $c = 2a = -2$: $y_p = -x^2 - 2$. General
solution $y = -x^2 - 2 + \lambda\eu^x + \mu\eu^{-x}$.

Cauchy: $y(0) = -2 + \lambda + \mu = 0$ and $y'(0) = \lambda - \mu =
1$: $\lambda = \frac32$, $\mu = \frac12$. So
$y = -x^2 - 2 + \frac{3\eu^x + \eu^{-x}}{2}$.
\end{solution}

\begin{solution}{exo:b1:diffeq:5}
$a(x) = \tan x$, $A(x) = -\ln(\cos x)$ (valid: $\cos > 0$ on the
interval), $\eu^{-A} = \cos x$: homogeneous solutions $\lambda\cos x$.
Variation of constants: $\mu'(x) = \sin 2x \cdot \frac{1}{\cos x} =
2\sin x$, so $\mu = -2\cos x + \lambda$ and
\[
y(x) = -2\cos^2 x + \lambda \cos x .
\]
\emph{Check:} $y' = 4\cos x \sin x - \lambda\sin x$ and $y\tan x =
-2\cos x\sin x + \lambda \sin x$; their sum is $2\cos x\sin x =
\sin 2x$, as required.
\end{solution}

\begin{solution}{exo:b1:diffeq:6}
$\chi(r) = r^2 - 4r + 3 = (r-1)(r-3)$: $\gamma = 1$ is a simple root
($m = 1$). Try $y_p = x(ax + b)\,\eu^x$. With $u = ax^2 + bx$,
\[
y_p'' - 4y_p' + 3y_p = \bigl(u'' + (2 - 4)u' + \chi(1) u\bigr)\eu^x
= \bigl(2a - 2(2ax + b)\bigr)\eu^x .
\]
Identify with $(2x + 1)\eu^x$: $-4a = 2$ and $2a - 2b = 1$, so $a =
-\frac12$, $b = -1$. General solution:
\[
y = -\Bigl(\frac{x^2}{2} + x\Bigr)\eu^{x} + \lambda\,\eu^{x} +
\mu\,\eu^{3x}, \qquad (\lambda, \mu) \in \R^2 .
\]
\end{solution}

\begin{solution}{exo:b1:diffeq:7}
Homogeneous: $y_h = \lambda\cos 2x + \mu\sin 2x$.

Right-hand side $x$ ($\gamma = 0$ not a root): $y_1 = ax + b$ with
$4(ax + b) = x$: $y_1 = \frac x4$.

Right-hand side $\sin 2x = \Im(\eu^{2\iu x})$, $2\iu$ simple root of
$r^2 + 4$: try $z = c\,x\,\eu^{2\iu x}$; then $z'' + 4z =
4\iu c\,\eu^{2\iu x}$, equal to $\eu^{2\iu x}$ for $c =
\frac{1}{4\iu} = -\frac{\iu}{4}$. So $z = -\frac{\iu x}{4}(\cos 2x +
\iu \sin 2x)$ and $y_2 = \Im(z) = -\frac{x\cos 2x}{4}$.

By superposition:
\[
y = \frac{x}{4} - \frac{x\cos 2x}{4} + \lambda\cos 2x + \mu\sin 2x .
\]
\end{solution}

\begin{solution}{exo:b1:diffeq:8}
The equation $T' + kT = 20k$ has constant particular solution $20$ and
homogeneous solutions $\lambda\eu^{-kt}$: $T(t) = 20 +
\lambda\,\eu^{-kt}$, and $T(0) = 80$ gives $\lambda = 60$:
\[
T(t) = 20 + 60\,\eu^{-kt} .
\]
$T(10) = 50$: $\eu^{-10k} = \frac12$, so $k = \frac{\ln 2}{10}$. Then
$T(t) = 25$ requires $\eu^{-kt} = \frac{5}{60} = \frac{1}{12}$, i.e.
\[
t = \frac{\ln 12}{k} = 10\,\frac{\ln 12}{\ln 2} \approx 35.8
\text{ minutes.}
\]
\end{solution}

\begin{solution}{exo:b1:diffeq:9}
\begin{enumerate}
  \item $\chi(r) = r^2 + 2\varepsilon r + 1$, $\Delta = 4(\varepsilon^2
        - 1)$.
        For $\varepsilon \in \intco{0}{1}$: roots $-\varepsilon \pm
        \iu\sqrt{1 - \varepsilon^2}$, so
        $y = \eu^{-\varepsilon t}\bigl(\lambda\cos\omega t +
        \mu\sin\omega t\bigr)$ with $\omega = \sqrt{1 -
        \varepsilon^2}$.
        For $\varepsilon = 1$: double root $-1$, $y = (\lambda + \mu
        t)\,\eu^{-t}$.
        For $\varepsilon > 1$: real roots $r_\pm = -\varepsilon \pm
        \sqrt{\varepsilon^2 - 1}$, both $< 0$, and $y =
        \lambda\eu^{r_+t} + \mu\eu^{r_-t}$.
  \item For $\varepsilon \in \intoo{0}{1}$: $\abs y \leq
        \eu^{-\varepsilon t}(\abs\lambda + \abs\mu) \to 0$. For
        $\varepsilon = 1$: $(\lambda + \mu t)\eu^{-t} \to 0$
        (exponential beats polynomial,
        \cref{prop:b1:functions:powerrules}). For $\varepsilon > 1$:
        both exponentials decay since $r_\pm < 0$ (indeed
        $\sqrt{\varepsilon^2 - 1} < \varepsilon$). For $\varepsilon =
        0$: $y = \lambda\cos t + \mu\sin t$ has constant amplitude
        $\sqrt{\lambda^2 + \mu^2} \neq 0$ unless $y = 0$.
  \item Write $\lambda\cos\omega t + \mu\sin\omega t =
        R\cos(\omega t - \varphi)$ with $R = \sqrt{\lambda^2 + \mu^2}
        > 0$. The zeros of $y$ are those of $\cos(\omega t -
        \varphi)$ (the factor $\eu^{-\varepsilon t}$ never
        vanishes): $\omega t - \varphi \equiv \frac\pi2 \pmod \pi$,
        an arithmetic progression with gap $\frac{\pi}{\omega} =
        \frac{\pi}{\sqrt{1 - \varepsilon^2}}$.
\end{enumerate}
\end{solution}

\begin{solution}{exo:b1:diffeq:10}
Set $x = y = 0$: $2f(0) = 2f(0)^2$, and $f(0) \neq 0$ gives $f(0) =
1$. Fix $x$ and differentiate the equation twice with respect to $y$:
\[
f''(x+y) + f''(x-y) = 2 f(x) f''(y) .
\]
Setting $y = 0$: $\;2f''(x) = 2 f(x) f''(0)$, that is
\[
f''(x) = c\,f(x), \qquad c = f''(0).
\]
\emph{Case $c = \omega^2 > 0$:} $f(x) = \lambda\cosh\omega x +
\mu\sinh\omega x$; $f(0) = 1$ gives $\lambda = 1$. Plugging into the
functional equation and using the addition formulas
(\cref{prop:b1:functions:hyprules}), the equation forces $\mu = 0$
(compare the coefficients of $\sinh\omega x \sinh\omega y$ or
evaluate at $x = y$): $f = \cosh\omega x$, which does satisfy
$\cosh(x+y) + \cosh(x-y) = 2\cosh x\cosh y$.

\emph{Case $c = -\omega^2 < 0$:} similarly $f(x) = \cos\omega x$
($\omega \neq 0$), which satisfies the equation.

\emph{Case $c = 0$:} $f$ affine with $f(0) = 1$: $f(x) = 1 + \mu x$;
the equation forces $\mu = 0$, excluded ($f$ not constant).

Conclusion: the solutions are $f(x) = \cos\omega x$ and $f(x) =
\cosh\omega x$, $\omega > 0$.
\end{solution}

\begin{solution}{exo:b1:diffeq:11}
Set $z(t) = y(\eu^t)$, so that $y(x) = z(\ln x)$ for $x > 0$. Then
\[
y'(x) = \frac{z'(\ln x)}x,
\qquad
y''(x) = \frac{z''(\ln x) - z'(\ln x)}{x^2} ,
\]
and substituting into the equation:
\[
x^2y'' - xy' + y = \bigl(z'' - z'\bigr) - z' + z
= z'' - 2z' + z = 0 .
\]
Characteristic polynomial $(r - 1)^2$: double root $1$, so $z(t) =
(\lambda + \mu t)\,\eu^t$ and, back in the variable $x = \eu^t$:
\[
y(x) = (\lambda + \mu\ln x)\,x,
\qquad \lambda, \mu \in \R .
\]
\end{solution}

\begin{solution}{exo:b1:diffeq:12}
\begin{enumerate}
  \item In normalized form $y' - \frac2x\,y = 0$ on each interval:
        $A(x) = -2\ln\abs x$, so the solutions are $y = a x^2$ on
        $\intoo0{+\infty}$ and $y = b x^2$ on $\intoo{-\infty}0$,
        with independent constants
        (\cref{thm:b1:diffeq:homogeneous1}).
  \item Let $y = ax^2$ for $x \geq 0$ and $bx^2$ for $x < 0$. On
        each open half-line $y$ is differentiable with $xy' = 2y$.
        At $0$: the difference quotients $\frac{y(h) - y(0)}h =
        ah$ or $bh$ tend to $0$, so $y'(0) = 0$ exists, and the
        equation at $x = 0$ reads $0 \cdot y'(0) = 2y(0) = 0$:
        satisfied. So $y$ solves the equation on all of $\R$.
  \item The solutions on $\R$ are exactly these glued functions: a
        \emph{two}-parameter family for a first-order equation.
        There is no contradiction with \cref{thm:b1:diffeq:voc},
        whose hypotheses fail here: written as $y' + a(x)y = 0$,
        the coefficient $a(x) = -\frac2x$ is not continuous at $0$
        --- indeed not defined --- so $\R$ is not an interval on
        which the theorem applies. The singularity at $0$
        disconnects the two half-lines, and the value $y(0) = 0$
        is forced, carrying no information across. Every Cauchy
        datum at $x_0 \neq 0$ determines the solution only on the
        half-line containing $x_0$.
\end{enumerate}
\end{solution}

\begin{solution}{pb:b1:diffeq:1}
\textbf{1.} $\chi(r) = r^2 + 2\lambda r + \omega_0^2$, $\Delta =
4(\lambda^2 - \omega_0^2) < 0$ for $0 < \lambda < \omega_0$: roots
$-\lambda \pm \iu\omega_d$, so by
\cref{thm:b1:diffeq:homogeneous2}
\[
x(t) = \eu^{-\lambda t}\bigl(\lambda_1\cos\omega_d t +
\mu_1\sin\omega_d t\bigr), \qquad (\lambda_1, \mu_1) \in \R^2 .
\]
For $\lambda = 0$: $x = \lambda_1\cos\omega_0 t +
\mu_1\sin\omega_0 t$. The critical ($\lambda = \omega_0$) and
overdamped ($\lambda > \omega_0$) regimes are those of
\cref{exo:b1:diffeq:9} (after rescaling time): $(\lambda_1 +
\mu_1 t)\eu^{-\lambda t}$, resp.\ combinations of $\eu^{r_\pm t}$
with $r_\pm = -\lambda \pm \sqrt{\lambda^2 - \omega_0^2}$.

\textbf{2.} Underdamped: $\abs x \leq \eu^{-\lambda
t}(\abs{\lambda_1} + \abs{\mu_1}) \to 0$. Critical: $(\lambda_1 +
\mu_1 t)\eu^{-\lambda t} \to 0$ since exponentials beat
polynomials (\cref{prop:b1:functions:powerrules}). Overdamped:
$r_- < r_+ = -\lambda + \sqrt{\lambda^2 - \omega_0^2} < 0$
because $\sqrt{\lambda^2 - \omega_0^2} < \lambda$; both
exponentials decay.

\textbf{3.} Along a solution of $(H)$, using $x'' = -2\lambda x' -
\omega_0^2 x$:
\[
\mathcal E'(t) = x'x'' + \omega_0^2 x x'
= x'\bigl(-2\lambda x' - \omega_0^2 x\bigr) + \omega_0^2 xx'
= -2\lambda\,x'^2 \leq 0 .
\]
If $x(t_0) = x'(t_0) = 0$ then $\mathcal E(t_0) = 0$; $\mathcal E$
is nonnegative and nonincreasing, so $\mathcal E \equiv 0$ on
$\intco{t_0}{+\infty}$, forcing $x \equiv 0$ there; for $t \leq
t_0$, run the same argument on $\tilde x(t) = x(2t_0 - t)$, which
solves the equation with damping $-\lambda$ but still has
$\tilde{\mathcal E}(t_0) = 0$ and $\tilde{\mathcal E}' =
+2\lambda\tilde x'^2 \geq 0$ with $\tilde{\mathcal E} \geq 0$;
nonnegative, nondecreasing and zero at the right end of
$\intoc{-\infty}{t_0}$ means zero throughout. So $x \equiv 0$ on
$\R$ --- an energy proof of uniqueness, valid for all $\lambda
\geq 0$.

\textbf{4.} $x(t + T_d) = R\,\eu^{-\lambda t}\eu^{-\lambda
T_d}\cos(\omega_d t + 2\pi - \varphi) = \eu^{-\lambda T_d}x(t)$.
The shrink factor per pseudo-period is $\eu^{-\delta}$ with
$\delta = \lambda T_d = \frac{2\pi\lambda}{\omega_d}$. For
$\omega_0 = 1$, $\lambda = 0.1$: $\omega_d = \sqrt{0.99} =
0.99499$, so $\delta = \frac{0.62832}{0.99499} = 0.6315$: each
swing keeps $\eu^{-0.63} \approx 53\%$ of its amplitude.

\textbf{5.} The amplitude factor is $\eu^{-\lambda t}$, which
decays by $\eu$ over $t = \frac1\lambda$. That interval contains
$\frac{1/\lambda}{T_d} = \frac{\omega_d}{2\pi\lambda}$
pseudo-periods. For $\lambda \ll \omega_0$, $\omega_d \approx
\omega_0$ and this is $\approx \frac{\omega_0}{2\pi\lambda} =
\frac Q\pi$. A $Q = 300$ guitar string rings for about a hundred
periods; a $Q = 1$ door damper does not complete one.

\textbf{6.} Substituting $z\,\eu^{\iu\Omega t}$ into the left-hand
side gives $z\,(-\Omega^2 + 2\iu\lambda\Omega +
\omega_0^2)\,\eu^{\iu\Omega t}$, which equals $A\,\eu^{\iu\Omega
t}$ exactly for $z = \frac{A}{\omega_0^2 - \Omega^2 +
2\iu\lambda\Omega}$ (the denominator is nonzero: its imaginary
part is $2\lambda\Omega > 0$). Since the coefficients are real,
the real part $x_p = \Re\bigl(z\eu^{\iu\Omega t}\bigr)$ solves the
equation with right-hand side $\Re\bigl(A\eu^{\iu\Omega t}\bigr) =
A\cos\Omega t$.

\textbf{7.} Write $\omega_0^2 - \Omega^2 + 2\iu\lambda\Omega =
\sqrt{D}\,\eu^{\iu\varphi}$ with $D = (\omega_0^2 - \Omega^2)^2 +
4\lambda^2\Omega^2$ and $\varphi \in \intoo0\pi$ (the imaginary
part $2\lambda\Omega$ is positive), so $\tan\varphi =
\frac{2\lambda\Omega}{\omega_0^2 - \Omega^2}$. Then $z =
\frac{A}{\sqrt D}\eu^{-\iu\varphi}$ and
\[
x_p(t) = \Re\Bigl(\frac A{\sqrt D}\,\eu^{\iu(\Omega t -
\varphi)}\Bigr) = R\cos(\Omega t - \varphi),
\qquad R = \frac A{\sqrt D} .
\]

\textbf{8.} As $\Omega \to 0^+$: $D \to \omega_0^4$, so $R \to
A/\omega_0^2$ and $\tan\varphi \to 0^+$ with $\varphi \in
\intoo0{\frac\pi2}$: $\varphi \to 0$. The mass follows the force
quasi-statically, displaced by force/stiffness. As $\Omega \to
+\infty$: $D \sim \Omega^4$, so $R \sim A/\Omega^2 \to 0$, and
$\varphi \to \pi$ (the complex number $\omega_0^2 - \Omega^2 +
2\iu\lambda\Omega$ goes to the second quadrant with argument $\to
\pi$): the mass barely moves, and in opposition of phase ---
inertia dominates.

\textbf{9.} By \cref{thm:b1:diffeq:structure2}~(1), every solution
is $x = x_p + x_h$ with $x_h$ solving $(H)$; by question 2, $x_h(t)
\to 0$, so $x(t) - x_p(t) \to 0$: all solutions converge to the
same steady state. The initial conditions only shape the transient.

\textbf{10.} If $x$ is a periodic solution, $x - x_p = x_h$ is a
periodic solution of $(H)$ which tends to $0$ at $+\infty$; a
periodic function with limit $0$ is identically $0$ (its values on
one period repeat forever, so every value is a limit of a
subsequence tending to $0$). Hence $x = x_p$.

\textbf{11.} Here $\lambda = 1$, $\omega_0^2 = 2$, $\Omega = 1$,
$A = 1$: $z = \frac1{2 - 1 + 2\iu} = \frac{1 - 2\iu}5$, so
\[
x_p = \Re\Bigl(\frac{(1 - 2\iu)(\cos t + \iu\sin t)}5\Bigr)
= \frac{\cos t + 2\sin t}5 .
\]
Homogeneous: roots of $r^2 + 2r + 2$ are $-1 \pm \iu$: $x_h =
\eu^{-t}(C\cos t + S\sin t)$. Conditions: $x(0) = \frac15 + C = 0$
gives $C = -\frac15$; differentiating, $x'(0) = \frac25 - C + S =
0$ gives $S = C - \frac25 = -\frac35$. Hence
\[
x(t) = \underbrace{\frac{\cos t + 2\sin t}5}_{\text{steady}}
- \underbrace{\eu^{-t}\,\frac{\cos t + 3\sin
t}5}_{\text{transient}} ,
\]
the transient dying like $\eu^{-t}$.

\textbf{12.} Expanding, $g(u) = u^2 - 2(\omega_0^2 - 2\lambda^2)u
+ \omega_0^4 = (u - u_r)^2 + g(u_r)$ with $u_r = \omega_0^2 -
2\lambda^2$ and
\[
g(u_r) = \omega_0^4 - u_r^2 = (\omega_0^2 - u_r)(\omega_0^2 +
u_r) = 2\lambda^2\bigl(2\omega_0^2 - 2\lambda^2\bigr) =
4\lambda^2(\omega_0^2 - \lambda^2) .
\]
If $2\lambda^2 < \omega_0^2$, then $u_r > 0$ is an admissible
squared frequency: $g$ has a strict minimum there, so $R = A/\sqrt
g$ has a strict maximum at $\Omega_r = \sqrt{u_r} =
\sqrt{\omega_0^2 - 2\lambda^2}$, with $R_{\max} = A/\sqrt{g(u_r)}
= \frac{A}{2\lambda\sqrt{\omega_0^2 - \lambda^2}}$.

\textbf{13.} $R(0) = A/\omega_0^2$, so
\[
\frac{R_{\max}}{R(0)}
= \frac{\omega_0^2}{2\lambda\sqrt{\omega_0^2 - \lambda^2}}
= \frac{\omega_0}{2\lambda}\cdot
\frac{\omega_0}{\sqrt{\omega_0^2 - \lambda^2}}
= Q\,\Bigl(1 - \frac{\lambda^2}{\omega_0^2}\Bigr)^{-1/2}
\geq Q .
\]
For weak damping the correction factor is close to $1$: resonance
multiplies the static displacement by essentially $Q$.

\textbf{14.} $V(\Omega)^2 = \frac{A^2 u}{g(u)}$ with $u =
\Omega^2$. Its derivative has the sign of $g(u) - u\,g'(u) =
(\omega_0^2 - u)^2 + 4\lambda^2 u - u\bigl(2(u - \omega_0^2) +
4\lambda^2\bigr) = (\omega_0^2 - u)^2 + 2u(\omega_0^2 - u) =
(\omega_0^2 - u)(\omega_0^2 + u)$, positive for $u < \omega_0^2$
and negative beyond: strict maximum exactly at $\Omega =
\omega_0$, for every damping. There $\tan\varphi$ blows up with
$\varphi \in \intoo0\pi$: $\varphi = \frac\pi2$, and $x_p'(t) =
-R\Omega\sin(\Omega t - \frac\pi2) = R\Omega\cos(\Omega t)$ is
exactly in phase with the force: optimal power transfer.

\textbf{15.} $R = R_{\max}/\sqrt2 \iff g(u) = 2g(u_r) \iff (u -
u_r)^2 = g(u_r) = 4\lambda^2(\omega_0^2 - \lambda^2)$, giving
\[
u_\pm = u_r \pm 2\lambda\sqrt{\omega_0^2 - \lambda^2},
\qquad
\Omega_+^2 - \Omega_-^2 = 4\lambda\sqrt{\omega_0^2 - \lambda^2} .
\]
Then $\Omega_+ - \Omega_- = \frac{\Omega_+^2 -
\Omega_-^2}{\Omega_+ + \Omega_-}$, and for $\lambda \ll \omega_0$
both $\Omega_\pm \approx \omega_0$: $\Omega_+ - \Omega_- \approx
\frac{4\lambda\omega_0}{2\omega_0} = 2\lambda$, so
$\frac{\omega_0}{\Omega_+ - \Omega_-} \approx
\frac{\omega_0}{2\lambda} = Q$. Measuring a resonance peak's width
measures its quality factor.

\textbf{16.} $Q = 10$; $\Omega_r = \sqrt{1 - 2(0.05)^2} =
\sqrt{0.995} = 0.9975$; $R_{\max} = \frac1{2 \times 0.05
\sqrt{1 - 0.0025}} = \frac1{0.1 \times 0.99875} = 10.01$; static
response $R(0) = 1$; bandwidth $\approx 2\lambda = 0.1$. A tall
thin spike of height $\approx Q$ over a plateau of height $1$.

\textbf{17.} If $2\lambda^2 \geq \omega_0^2$ then $u_r \leq 0$ and
$g'(u) = 2(u - u_r) > 0$ for all $u > 0$: $g$ increases strictly
on $\intoo0{+\infty}$, so $R = A/\sqrt g$ decreases strictly from
$R(0) = A/\omega_0^2$: the response is largest at zero frequency
and there is no peak.

\textbf{18.} $\gamma = \iu\Omega$ is not a root of $r^2 +
\omega_0^2$ (as $\Omega \neq \omega_0$), so
\cref{met:b1:diffeq:particular} with $m = 0$ gives $x_p =
\frac{A\cos\Omega t}{\omega_0^2 - \Omega^2}$ (substitute and
check: $-\Omega^2 + \omega_0^2$ times the cosine). General
solution:
\[
x(t) = \frac{A\cos\Omega t}{\omega_0^2 - \Omega^2}
+ \lambda_1\cos\omega_0 t + \mu_1\sin\omega_0 t .
\]

\textbf{19.} $x(0) = 0$ forces $\lambda_1 = -\frac A{\omega_0^2 -
\Omega^2}$, and $x'(0) = 0$ forces $\mu_1 = 0$:
\[
x(t) = \frac{A}{\omega_0^2 - \Omega^2}\,
\bigl(\cos\Omega t - \cos\omega_0 t\bigr)
= \frac{2A}{\omega_0^2 - \Omega^2}\,
\sin\Bigl(\frac{(\omega_0 - \Omega)t}2\Bigr)
\sin\Bigl(\frac{(\omega_0 + \Omega)t}2\Bigr) ,
\]
by the product formula $\cos a - \cos b =
2\sin\frac{b - a}2\sin\frac{b + a}2$ applied with $a = \Omega t$,
$b = \omega_0 t$.

\textbf{20.} For $\Omega$ near $\omega_0$, the second sine
oscillates at the fast frequency $\frac{\omega_0 + \Omega}2
\approx \omega_0$, while the first is a slow envelope of frequency
$\frac{\abs{\omega_0 - \Omega}}2$: the amplitude of the fast
oscillation waxes and wanes with envelope period
$\frac{2\pi}{\abs{\omega_0 - \Omega}}$ (two beats per envelope
period), reaching maxima $\frac{2A}{\abs{\omega_0^2 - \Omega^2}}$.
As $\Omega \to \omega_0$, the beats become both slower (period
$\to \infty$) and taller (amplitude $\to \infty$).

\textbf{21.} Fix $t$. As $\Omega \to \omega_0$:
\[
\frac{2A}{\omega_0^2 - \Omega^2}\sin\Bigl(\frac{(\omega_0 -
\Omega)t}2\Bigr)
= \frac{2A}{\omega_0 + \Omega}\cdot
\frac{\sin\bigl(\frac{(\omega_0 - \Omega)t}2\bigr)}
{\omega_0 - \Omega}
\longrightarrow \frac{2A}{2\omega_0}\cdot\frac t2
= \frac{At}{2\omega_0} ,
\]
while $\sin\bigl(\frac{(\omega_0 + \Omega)t}2\bigr) \to
\sin(\omega_0 t)$: the limit is $x_\infty(t) =
\frac{At\sin\omega_0 t}{2\omega_0}$. Direct check: with $C =
\frac A{2\omega_0}$, $x_\infty = Ct\sin\omega_0 t$ has
$x_\infty'' = 2C\omega_0\cos\omega_0 t - C\omega_0^2 t
\sin\omega_0 t$, so $x_\infty'' + \omega_0^2 x_\infty =
2C\omega_0\cos\omega_0 t = A\cos\omega_0 t$, with $x_\infty(0) =
0$ and $x_\infty'(0) = C\sin 0 + C\omega_0 \cdot 0 \cdot \cos 0 =
0$ --- for $\omega_0 = A = 1$ this is
exactly \cref{ex:b1:diffeq:oscillation}. Resonance is the
degeneration of beats: the first swell of the envelope, stretched
to infinite length.

\textbf{22.} Without damping the resonant amplitude grows linearly
and without bound; with damping $\lambda > 0$ the growth saturates
at $R_{\max} \approx Q\,\frac A{\omega_0^2}$. The mechanism:
dissipation removes energy at a rate growing with the amplitude
(question 23), so the build-up stops exactly when the damping
burns energy as fast as the forcing supplies it.

\textbf{23.} With $x_p' = -R\Omega\sin(\Omega t - \varphi)$, over
a period the averages $\langle\cos^2\rangle =
\langle\sin^2\rangle = \frac12$ and $\langle\sin\cos\rangle = 0$
give:
\[
\langle P_{\mathrm{diss}}\rangle
= 2\lambda\,R^2\Omega^2\,\langle\sin^2(\Omega t -
\varphi)\rangle = \lambda R^2\Omega^2 ;
\]
and expanding $\sin(\Omega t - \varphi) = \sin\Omega t\cos\varphi
- \cos\Omega t\sin\varphi$:
\[
\langle P_{\mathrm{in}}\rangle
= -AR\Omega\,\bigl\langle\cos\Omega t\,\sin(\Omega t -
\varphi)\bigr\rangle
= AR\Omega\,\frac{\sin\varphi}2 .
\]
Since $\sin\varphi = \frac{2\lambda\Omega}{\sqrt D} =
\frac{2\lambda\Omega R}A$, this is $\frac{AR\Omega}2 \cdot
\frac{2\lambda\Omega R}A = \lambda R^2\Omega^2$: injected and
dissipated powers balance exactly --- the defining property of a
steady regime.

\textbf{24.} (i) The structure theorem split every solution into
steady state plus transient (questions 9--11) and reduced
uniqueness to the homogeneous problem. (ii) The complex method
turned the search for a particular solution into one division of
complex numbers (question 6), with amplitude and phase read off a
modulus and an argument. (iii) The resonance curve is a pure
function study --- a quadratic in $u = \Omega^2$, its minimum,
its level sets --- in the style of \cref{ch:b1:functions}
(questions 12--17).

\textbf{25.} Free and damped: decaying pseudo-oscillations,
lifetime $\frac1\lambda$, about $\frac Q\pi$ swings. Forced and
damped: transients die, and a unique sinusoidal steady state
survives at the forcing frequency, with amplitude peaking near
$\omega_0$ (height $\approx Q \times$ static, width $\approx
\frac{\omega_0}Q$) and phase sweeping from $0$ to $\pi$ through
$\frac\pi2$ at $\omega_0$. Free and undamped: perpetual
oscillation. Forced and undamped: beats, degenerating into
linearly growing resonance at exact tuning. One dimensionless
number, $Q = \frac{\omega_0}{2\lambda}$, tunes everything ---
peak height, bandwidth, and transient lifetime are three readings
of the same dial. The sequel: rewriting $x'' + 2\lambda x' +
\omega_0^2x$ as a first-order system opens the matrix methods of
\cref{ch:b1:matrices} and the Year 2 volume, and decomposing an
arbitrary periodic forcing into sinusoids (Fourier series, Year 3
volume) makes this problem's single-frequency analysis the
universal building block: solve for each frequency, superpose.
\end{solution}
